\documentclass[11pt]{article}
\usepackage[T1]{fontenc}
\usepackage{lmodern}
\usepackage[margin=1in,top=0.9in,bottom=0.9in]{geometry}
\usepackage{amsmath,amssymb,amsthm,mathtools}
\usepackage{microtype,booktabs,longtable,array}
\usepackage{xcolor,enumitem,listings,fancyhdr}
\definecolor{ink}{RGB}{25,52,78}
\definecolor{papergray}{RGB}{246,248,250}
\usepackage[colorlinks=true,linkcolor=ink,urlcolor=ink,citecolor=ink]{hyperref}
\hypersetup{pdftitle={MAIS-O11: Consolidated Progress and Lean Certificates},
  pdfauthor={Research record for Andrew Hagy},
  pdfsubject={Proof lengths, Lob's theorem, PA-bin, and verification scope}}
\setlength{\headheight}{14pt}
\pagestyle{fancy}\fancyhf{}
\fancyhead[L]{\small\textcolor{ink}{MAIS-O11 \;|\; Consolidated progress}}
\fancyhead[R]{\small 23 September 2026}
\fancyfoot[C]{\small\thepage}
\renewcommand{\headrulewidth}{0.3pt}
\setlength{\emergencystretch}{2em}
\setlist{itemsep=3pt,topsep=5pt}
\lstset{basicstyle=\small\ttfamily,breaklines=true,columns=fullflexible,
  keepspaces=true,backgroundcolor=\color{papergray},frame=single,
  rulecolor=\color{black!15},xleftmargin=5pt,xrightmargin=5pt,
  aboveskip=9pt,belowskip=9pt,showstringspaces=false}
\newtheorem{theorem}{Theorem}[section]
\newtheorem{proposition}[theorem]{Proposition}
\newtheorem{lemma}[theorem]{Lemma}
\newtheorem{corollary}[theorem]{Corollary}
\theoremstyle{definition}
\newtheorem{definition}[theorem]{Definition}
\theoremstyle{remark}
\newtheorem{remark}[theorem]{Remark}
\newcommand{\PA}{\mathsf{PA}}
\newcommand{\PAbin}{\mathsf{PA}_{\mathrm{bin}}}
\newcommand{\Con}{\operatorname{Con}}
\newcommand{\Prf}{\operatorname{Prf}}
\newcommand{\Prov}{\operatorname{Prov}}
\newcommand{\gn}[1]{\ulcorner #1\urcorner}
\newcommand{\num}[1]{\overline{#1}}
\newcommand{\N}{\mathbb N}
\newcommand{\cO}{\mathcal O}
\newcommand{\status}[1]{\par\smallskip\noindent\textbf{Verification scope.} #1\par\smallskip}
\title{\textcolor{ink}{Does L\"ob's rule shorten proofs?}\\[5pt]
  \large MAIS-O11: consolidated progress and Lean certificates}
\author{Research record for Andrew Hagy\\
  \small AI-assisted mathematical development and verification audit}
\date{23 September 2026}
\begin{document}
\maketitle
\begin{abstract}
This report consolidates the slow-verifier and connective-tag constructions,
the conditional quantitative L\"ob conversion, the raw-string quotation and
trace compiler, the proof-assembly obstruction, and the latest finite-budget
envelope theorem. The strongest new structural conclusion is exact: for a
fixed provability predicate and toll, absence of a part~1 witness is equivalent
to sublinearity of the maximum excess over the reflection-proof length. This
negative condition implies a joint linear bound for the overhead function.
Conversely, failure of every joint linear bound already yields witnesses with
every fixed integer factor. These implications are verified in Lean.

Neither part has been decided for the intended ordinary $\PAbin$ checker.
The supplied certificates establish the precise abstract, syntactic, and
numerical statements identified below; they do not instantiate every
arithmetical hypothesis of the earlier chosen-presentation constructions.
The accompanying cumulative archive preserves six Lean projects, their
research notes, and a single command for rebuilding and auditing all six.
\end{abstract}

\paragraph{How to read this report.}
Sections~\ref{sec:envelope}--\ref{sec:consequences} give the latest structural
result and its full mathematical proof. Sections~\ref{sec:early}--\ref{sec:obstruction}
record the earlier mechanisms, implementation milestones, and failed transfer
routes. Section~\ref{sec:ledger} maps the claims to certificates and explains
how to verify them locally. The report supersedes the earlier manuscript as
the summary of \emph{verification status}; the earlier manuscript and detailed
milestone notes remain in the archive as research history.

\clearpage
{\small\tableofcontents}
\clearpage

\section{The fixed question and the scope of the comparison}
\label{sec:question}

Fix one proof system $S$, one arithmetization $\Box_S$ of its checker, and
one nonnegative integer $c=C_1$ for the cheap-reflection toll. For a provable
sentence $P$, write
\begin{equation}\label{eq:notation}
 d(P)=\ell_S(P),\qquad r(P)=\ell_S(\Box_S P\to P),\qquad
 s(P)=|P|,\qquad t(P)=c(s(P)+1).
\end{equation}
The agenda's cheap direction is $r(P)\le d(P)+t(P)$.
Part~1 asks for a fixed $\delta>0$ and a sequence $P_i$ such that
\begin{equation}\label{eq:partone}
 r(P_i)\longrightarrow\infty,\qquad
 d(P_i)\ge(1+\delta)r(P_i)+t(P_i).
\end{equation}
Part~2 concerns the fixed ordinary $\PAbin$ presentation and the function
\begin{equation}\label{eq:F}
 F_S(k,n)=\max\bigl(\{d(P):r(P)\le k,\ s(P)\le n\}\cup\{0\}\bigr).
\end{equation}
It asks whether this function has a fixed polynomial upper bound in $k+n$.
Adding $1$ to this argument is harmless for growth; the ordinary checker
accepts no empty proof, so the zero-budget corner can also be handled directly.
The maximum includes $0$ to cover an empty set of eligible sentences.

\subsection{A necessary finite-proof convention}
The agenda allows $\ell_S(P)=\infty$ for unprovable $P$. If the displayed
part~1 condition is read literally in the extended naturals, without requiring
finite proofs, it has an unintended loophole. In a consistent system satisfying
L\"ob's theorem, both $\bot$ and $\Box_S\bot\to\bot$ are unprovable.
The constant sequence $P_i=\bot$ then has $d(P_i)=r(P_i)=\infty$ and
satisfies an extended-arithmetic reading of the inequality and divergence.
This is not proof shortening. Throughout this report, part~1 concerns
\emph{provable} sentences and finite minimum lengths. L\"ob's theorem ensures
that a sentence with a reflection proof has a direct proof; the cheap direction
gives the converse for the families under consideration.

\subsection{What is being compared}
A shortest completed $S$-proof of $P$ may itself use a formalized L\"ob argument.
Thus $d(P)$ already minimizes over proofs obtained that way. The quantity
$r(P)$ measures the premise \emph{before} the cost of converting it into an
$S$-proof of $P$ has been paid. One can equivalently view the reflection file
as a new certificate format with an external L\"ob-rule step about the old box.
Changing $\Box_S$ changes the question even when the two checkers accept the
same standard files.

The headline question about saving one symbol is weaker than
\eqref{eq:partone}. Unbounded additive savings alone need not produce a fixed
fractional saving after the toll. For example, the numerical pattern
$d=r+\lfloor\sqrt r\rfloor$, with zero toll, has unbounded additive excess
but no fixed-factor family. This is an illustration of the quantifiers, not
an asserted family of PA sentences.

\subsection{The intended PA-bin presentation}
The uploaded MAIS-A1 fixes a Hilbert presentation based on Enderton, binary
numerals, a finite character alphabet, tagged proof records, and priced nested
abbreviations. Definitions may name \emph{arbitrary strings}; later fragments
may reference earlier definitions. The checker expands these definitions and
checks the resulting proof. Hence a short source file may expand enormously.
The specified logical axiom class includes all propositional tautologies,
which must not silently be replaced by a fixed finite Frege basis.

Two distinctions matter for the remaining work. A polynomial bound in expanded
length need not be polynomial in written file length. Also, a Lean proof that
a compiler preserves values is not an object-level PA proof of the fixed
arithmetized checker predicate. These issues motivate the quotation work in
Section~\ref{sec:quotation}.

\section{The finite-budget envelope: an exact characterization}
\label{sec:envelope}

This section assumes only finite-valued $d,r,s$ on the target set and that
there are finitely many targets with $r(P)\le k$ for each $k$. In a proof system
over a finite alphabet, that finiteness follows because there are finitely
many files of bounded length and each file has a single conclusion. It does
not require a bound on the length of an expanded conclusion.

\begin{definition}[Excess and envelope]
Set
\begin{equation}\label{eq:envelope}
 e(P)=\max\{d(P)-r(P)-t(P),0\},\qquad
 H(k)=\max\bigl(\{e(P):r(P)\le k\}\cup\{0\}\bigr).
\end{equation}
The envelope $H$ is nondecreasing and finite, and
\begin{equation}\label{eq:universal-envelope}
 d(P)\le r(P)+t(P)+H(r(P)).
\end{equation}
A natural-valued function $h$ is called sublinear when
\begin{equation}\label{eq:sublinear}
 \forall q\in\N_{>0}\ \exists K\ \forall k\ge K:\quad qh(k)<k.
\end{equation}
This is the usual statement $h(k)=o(k)$, written without real division.
\end{definition}

It suffices to consider factors $1+1/q$: for any $\delta>0$, choose an
integer $q\ge1$ with $1/q\le\delta$. For $r(P)>0$, the separation with
this factor is equivalent to
\begin{equation}\label{eq:separated}
 (q+1)r(P)+qt(P)\le qd(P)
 \quad\Longleftrightarrow\quad r(P)\le q e(P).
\end{equation}
Indeed, separation forces $d(P)>r(P)+t(P)$, after which the equivalence
follows by subtraction. If $d(P)\le r(P)+t(P)$, both sides are false.

\begin{theorem}[Exact negative condition]\label{thm:envelope}
The following conditions are equivalent.
\begin{enumerate}[label=(\roman*)]
\item There is no family satisfying \eqref{eq:partone}.
\item $H(k)=o(k)$.
\item There is a function $h=o(k)$ such that, for every target $P$,
\[
 d(P)\le r(P)+t(P)+h(r(P)).
\]
\end{enumerate}
\end{theorem}
\begin{proof}
First rewrite (i) as the cutoff condition
\begin{equation}\label{eq:cutoff}
 \forall q\ge1\ \exists K\ \forall P\quad
 r(P)\ge K\ \Longrightarrow\ q e(P)<r(P).
\end{equation}
If this condition fails for some $q$, there are witnesses at arbitrarily
large positive reflection lengths by \eqref{eq:separated}. Choose one with
$r(P_i)\ge i+1$ for each $i$. This gives the required divergent sequence.
Conversely any separating sequence eventually exceeds a proposed cutoff.
Using $K+1$ where necessary avoids the exceptional case $r(P)=0$.

Suppose \eqref{eq:cutoff} holds and fix $q$. Choose its cutoff $K$ and take
$k\ge qH(K)+1$. If $H(k)=0$, then $qH(k)<k$. Otherwise, choose an
attaining target $P$ with $r(P)\le k$ and $e(P)=H(k)$.
If $r(P)\ge K$, the cutoff gives
$qH(k)=qe(P)<r(P)\le k$. If $r(P)<K$, then $e(P)\le H(K)$ and
$qH(k)\le qH(K)<k$. Thus $H=o(k)$.

Conversely, if $H=o(k)$, choose a cutoff with $qH(k)<k$. For any target
whose reflection length exceeds that cutoff,
$qe(P)\le qH(r(P))<r(P)$, proving \eqref{eq:cutoff}.
This establishes (i)$\Leftrightarrow$(ii).

Condition (ii) implies (iii) by taking $h=H$ in
\eqref{eq:universal-envelope}. If (iii) holds, then $e(P)\le h(r(P))$.
The sublinearity cutoff for $h$ immediately gives \eqref{eq:cutoff}, proving
(iii)$\Rightarrow$(i).
\end{proof}

\status{The equivalences, cutoff formulation, and witness-sequence selection
are proved in \texttt{Envelope.lean}. Earlier notes treated a sublinear additive
converter bound as a stronger sufficient negative condition. The theorem
above improves that statement: in the finite-budget setting it is equivalent
to the negative answer. It does not prove that the condition holds for PA-bin.}

\subsection{Finite files supply the interface}
The module \texttt{FiniteFiles.lean} enumerates all words of length at most
$k$ over $\operatorname{Fin}(a)$, decodes them, and collects successful targets.
Its minimum-length premise is precisely
\begin{equation}\label{eq:minspec}
 r(P)\le k\ \Longleftrightarrow\
 \exists w\ (|w|\le k\ \wedge\ \operatorname{decode}(w)=\operatorname{some}(P)).
\end{equation}
The resulting finite target list supplies the hypothesis used in
Theorem~\ref{thm:envelope}. The decoder is a parameter, not an implemented PA
parser. Equation~\eqref{eq:minspec} is the defining property of minimum
reflection length; it is not a hidden polynomial-internalization assumption.

\subsection{Computability and proof-search time}
For a decidable checker with recognizable reflection conclusions, $H$ is
total computable as a metamathematical matter. Enumerate accepted reflection
files of length at most $k$, obtaining finitely many targets and their minimum
reflection lengths. For each target, search ordinary proofs in increasing
length until a direct proof is found. L\"ob's theorem guarantees termination.
Compute the maximum excess.

There is no useful runtime estimate in this argument. Similarly, shortest
direct-proof search converts an accepted reflection file to an ordinary proof;
if $H$ is sublinear, its output length satisfies the envelope bound. This does
not give a fast conversion algorithm. The current Lean project verifies the
finite-list and numerical results, not the computability argument or a PA
proof-search implementation.

\section{Consequences linking parts 1 and 2}
\label{sec:consequences}

\begin{theorem}[Near-unit global bounds]\label{thm:affine}
If there is no part~1 witness, then for each integer $q\ge1$ there is a
constant $B_q$ such that, for every target,
\begin{equation}\label{eq:affine}
 qd(P)\le(q+1)r(P)+qt(P)+B_q.
\end{equation}
In particular, for some constant $B$,
\begin{equation}\label{eq:Flinear}
 F_S(k,n)\le 2k+c(n+1)+B
 \le (2+c+B)(k+n+1).
\end{equation}
\end{theorem}
\begin{proof}
Choose $K$ from \eqref{eq:cutoff} and put $B_q=qH(K)$. If $r(P)\ge K$,
then $qe(P)<r(P)$. If $r(P)<K$, then $qe(P)\le qH(K)$. In both cases,
\[
 qe(P)\le r(P)+qH(K).
\]
Multiply $d(P)\le r(P)+t(P)+e(P)$ by $q$ and substitute this estimate to
obtain \eqref{eq:affine}. Set $q=1$ and maximize over $r(P)\le k$,
$s(P)\le n$ to get the first inequality in \eqref{eq:Flinear}; elementary
arithmetic gives the second.
\end{proof}

\begin{theorem}[Nonlinearity gives every fixed factor]\label{thm:nonlinear}
Suppose there is no constant $C$ such that
$F_S(k,n)\le C(k+n+1)$ for all $k,n$. Then, for every integer $a\ge1$ and
every cutoff $K$, there is a target $P$ with
\begin{equation}\label{eq:anyfactor}
 r(P)\ge K,\qquad d(P)\ge a r(P)+c(s(P)+1).
\end{equation}
For every $a\ge2$ this supplies a part~1 family with $\delta=a-1$.
\end{theorem}
\begin{proof}
Suppose \eqref{eq:anyfactor} fails for particular $a,K$. For all targets
with $r(P)\ge K$ we have $d(P)<ar(P)+t(P)$. For a target with $r(P)<K$,
\eqref{eq:universal-envelope} and $a\ge1$ give
\[
 d(P)\le r(P)+t(P)+H(K)\le ar(P)+t(P)+H(K).
\]
Hence this last bound holds for every target. Maximizing gives
\[
 F_S(k,n)\le ak+c(n+1)+H(K)
 \le (a+c+H(K))(k+n+1),
\]
contrary to the hypothesis. Select witnesses with $r(P_i)\ge i+1$ to obtain
divergence.
\end{proof}

These theorems refine the research alternatives. A negative answer to part~1
would imply the positive answer to part~2 in the stronger linear form.
A superpolynomial lower bound for part~2 would imply part~1; even failure of
all joint linear bounds suffices. The reverse implication is not established:
a fixed-factor gap can coexist with a linear bound having a larger constant.
Thus proving a polynomial upper bound alone would not rule out part~1.

\status{The integer statements and implications in this section are proved
in \texttt{Envelope.lean}. Their hypotheses are not proved for the actual
PA-bin length functions. The result narrows the alternatives without selecting
one of them.}

\section{Earlier chosen-presentation constructions}
\label{sec:early}

The first attempts deliberately used the freedom to choose $S$ in part~1.
Their essential mechanism was a slow finite test built into the checker.
This section records the mathematical constructions and their limitations;
the archive retains the complete earlier manuscript and tagged-checker note.

\subsection{Finite diagonal hardness and the first-pass certificate}
For a consistent ordinary base calculus $B$, a parameterized diagonal family
can express
\[
 D_n\ \longleftrightarrow\
 \text{there is no $B$-proof of $D_n$ of written length at most $n$}.
\]
If a short proof existed, numeralwise verification of that particular proof
would yield the negation of the displayed assertion, while the proof of
$D_n$ and the diagonal equivalence yield the assertion itself. Consistency
therefore gives $\ell_B(D_n)>n$. The bounded search actually terminates with
a positive result, and PA can verify that finite computation, so each
standard $D_n$ is provable. This says nothing favorable about the length of
that verification.

The first-pass project formalizes the abstract contradiction, L\"ob's rule
from explicit derivability and diagonal hypotheses, existence of shortest
proofs, and a quantitative transfer. In the transfer, a linear interpretation
of source proofs gives $n<C(d_n+1)$, while
\[
 r_n\le a(\lfloor\log_2 n\rfloor+2),\qquad
 s_n+1\le b(\lfloor\log_2 n\rfloor+2)
\]
supply a factor-two subsequence, including any fixed integer toll. Finitely
many bounded-length files and distinct reflection conclusions prove that
the minimum reflection lengths diverge.

\status{The theorem \texttt{SeparationHypotheses.conditional\_main} is
conditional. No value of its complete hypothesis structure for PA or the
guarded verifier has been constructed. Its axiom audit does not remove those
hypotheses.}

\subsection{The nullary-predicate slow verifier}
The earlier manuscript adds primitive nullary predicates $p_n$, whose binary
names have logarithmic length. It allows the reflection axiom
$\Box_Sp_n\to p_n$ and checks, whenever $p_n$ occurs, the finite condition
that $D_n$ has no short $B$-proof. A parameterized diagonal construction of
the checker formula aligns $\Box_S$ with this actual checking procedure,
avoiding an informal circular definition.

The interpretation $p_n\mapsto D_n$ takes the reflection axioms to instances
of a fixed base-theory theorem supplied by the guard. With whole-expression
abbreviations and the specified proof-record convention, the manuscript gives
a linear written-length translation. It derives, under consistency of $B$,
\[
 |p_n|+1=\cO(\log(n+2)),\qquad
 r_S(p_n)=\cO(\log(n+2)),\qquad d_S(p_n)>n/C-1.
\]
This is a proposed affirmative construction for the literal chosen-system
clause, with a slow checker and an expanded signature. It is not the fixed
PA-bin presentation. The actual checker arithmetization and translation have
not been implemented end to end in Lean.

An arithmetic-only variant used a recognizable family $G_n$ and a guard
depending on the numeral instance. Although its accepted standard files agree
with the ordinary checker's files under consistency, it fails a stronger
uniform instantiation property: for $\tau(x)=G(x)\to G(x)$, a uniform theorem
asserting $\Box_*\tau(\num n)$ would imply all the hard guards at once.
That variant therefore cannot supply the agenda's additional uniform bounded
principles. This failure motivated connective tags.

\subsection{Connective tags repair the substitution defect}
Choose an ordinary finite-basis PA calculus $B$ with primitive conjunction
and whole-expression abbreviations. Let
\[
 T_0=(0=0),\qquad T_{n+1}=(0=0)\mathbin{\&}T_n,
 \qquad P_n=T_n\mathbin{\&}D(\num n),
\]
where the diagonal construction is performed for the tagged formula itself,
and $D(n)$ asserts that $P_n$ has no $B$-proof of length at most $2^n$.
Guard indices are read from conjunction shapes, ignoring the terms in atomic
equalities. Capture-avoiding term substitution preserves these shapes and
cannot manufacture a new guard. MP, generalization, and universal
instantiation preserve the required guard conditions.

The guarded checker $S$ runs the ordinary check and all the finite guard tests.
Consistency gives truth of every standard guard, so the accepted standard
files, conclusions, and written lengths agree with $B$. Their arithmetic
provability predicates nevertheless differ. The guard supplies the implication
$\Box_SP_n\to D(\num n)$; adjoining a polynomial-size proof of $T_n$
gives the mathematical estimates
\begin{equation}\label{eq:tagged}
 d_S(P_n)>2^n,\qquad r_S(P_n)\le a(n+1)^d,\qquad
 s(P_n)+1\le b(n+1).
\end{equation}
These estimates imply a factor-two subsequence and superpolynomial $F_S$.
Their proving theory and box are both the chosen guarded $S$.

For the ordinary box, the same family has short proofs of $\Box_BP_n$.
Consequently an ordinary reflection proof, or a bridge proof
$\Box_BP_n\to\Box_SP_n$, would complete a direct proof after polynomial
additional work. The construction therefore also supplies the mathematical
bridge obstruction
\[
 \ell_B(\Box_BP_n\to P_n)+q_1(n)>2^n,\qquad
 \ell_B(\Box_BP_n\to\Box_SP_n)+q_2(n)>2^n
\]
for suitable polynomials $q_1,q_2$ under its stated proof-cost conventions.
Identity of standard proof files does not furnish cheap internal agreement
between the boxes.

\status{\texttt{TaggedSyntax.lean} proves concrete syntax, substitution,
guard, and expanded-trace invariants. \texttt{TaggedBounds.lean} proves the
numerical consequences of estimates such as \eqref{eq:tagged}. The ordinary
checker, diagonal arithmetization, guard truth, and the PA proof-length
estimates still have unformalized components. This is a mathematical
chosen-presentation construction with partial certificates, not a complete
Lean solution of part~1 or an answer about $F_{\PAbin}$.}

\section{The conditional polynomial upper-bound route}
\label{sec:upper}

Put $R_A=(\Box A\to A)$ and suppose a proof of $R_A$ has length $k$, with
$|A|=n$. A formalized L\"ob conversion has the following form:
\[
 \text{proof of }R_A
 \ \leadsto\ \text{proof of }\Box R_A
 \ \leadsto\ \text{proof of }\Box A
 \ \leadsto\ \text{proof of }A.
\]
The middle step uses the internal L\"ob theorem
$\Box(\Box A\to A)\to\Box A$; the last step uses the original $R_A$.

The audit project makes the required quantitative operations explicit.
Suppose, for fixed natural constants and exponents, that:
\begin{itemize}
\item internalizing the given proof costs at most $u(k+n+1)^d$;
\item the internal L\"ob instance has a proof of length at most $v(n+1)^e$;
\item MP combines two proofs with their summed lengths plus
      $m(|X|+|Y|+1)$;
\item boxed formulas and reflection formulas have sizes at most
      $b(|X|+1)$ and $r_0(|X|+1)$, respectively.
\end{itemize}
Define
\[
 W=m\bigl(b(r_0+3)+2\bigr),\qquad D=\max\{d,e,1\},\qquad
 K=1+u+v+W.
\]
Then the constructed proof has length at most
\begin{equation}\label{eq:polyupper}
 k+u(k+n+1)^d+v(n+1)^e+W(n+1)
 \le K(k+n+1)^D.
\end{equation}
If $d\le1$, the sharper bound is
\begin{equation}\label{eq:linear-premise}
 (u+1)k+(u+v+W)(n+1)^{\max\{e,1\}}.
\end{equation}

The logical internal L\"ob theorem is also derived from explicit HBL and
fixed-point inputs. For a fixed point $G\leftrightarrow(\Box G\to A)$,
the standard derivation gives $\Box G\to\Box A$, then $R_A\to G$, then
$\Box R_A\to\Box G\to\Box A$. This logical proof alone supplies no
polynomial character bound for its instances in the actual checker.

The first-pass logarithmic-reflection/linear-hardness estimates would rule out
every polynomial overhead bound: along $n=2^j$, their upper budgets grow
polynomially in $j$ while the direct lower bound is exponential. The audit
therefore proves that the first-pass separation hypotheses and a full
polynomial-conversion interface cannot both hold for the same presentation.
This identifies an incompatibility of assumptions, not a contradiction in PA.

\status{\texttt{AuditPass.lean} constructs the conversion and bounds its cost
from the explicit \path{QuantitativeLobTools} interface. It proves the
incompatibility just described.
There is no instance of \texttt{QuantitativeLobTools} for the complete actual
PA-bin checker. In particular, the first bullet above remains a substantive
obligation. Even \eqref{eq:linear-premise} generally has coefficient $u+1$
on $k$, and does not settle the near-unit question in part~1.}

\section{Raw-string quotation and compact arithmetic traces}
\label{sec:quotation}

The quotation project addresses the agenda's arbitrary-string abbreviation
grammar. Its input is a well-scoped parsed grammar: each definition is a list
of literal alphabet symbols and references to earlier definitions. Fragments
may cross term or formula boundaries. For $n$ definitions, set
\[
 M=n+\text{the total number of atoms in all right-hand sides}.
\]
Alphabet symbols are digits below a fixed base $b\ge2$. For an expanded word
$w$, its summary is
\[
 L_w=|w|,\qquad S_w=b^{|w|},\qquad V_w=\operatorname{raw}_b(w).
\]
The marked base-$b$ string code is $S_w+V_w$, which retains leading zeros
and distinguishes lengths.

\subsection{A compiler that avoids printing the expanded numeral}
For concatenation, the summaries obey
\begin{equation}\label{eq:concat}
 L_{xy}=L_x+L_y,\qquad S_{xy}=S_xS_y,\qquad
 V_{xy}=V_xS_y+V_y.
\end{equation}
The empty word has summary $(0,1,0)$; a literal digit $a$ has summary
$(1,b,a)$. The compiler uses three arithmetic registers per source definition
and generates terms containing only binary numerals, addition, multiplication,
and references to previous registers. Thus $S_w=b^{|w|}$ is a semantic
identity; an exponentiation symbol or the full huge numeral is not printed.

Lean proves, by induction through the grammar, that the generated terms have
the three intended values. It also proves identifier framing and injectivity,
freshness, and serialized character bounds. With
\[
 Q_b(M)=4(b+4)M^2+3M,
\]
the generated definitions use at most
\begin{equation}\label{eq:quote-bound}
 (6M+20)Q_b(M)
\end{equation}
characters. This is a bound in the parsed grammar mass, not expanded length.

\subsection{Local equality certificates}
The project adds ordinary closed-term equality proofs for the generated
definitions. After a definition $u:=t$, the expanded equality $u=t$ is $t=t$.
A shared proof of $\forall x\,x=x$, its universal-instantiation axiom, and MP
prove the local equality. A restricted Enderton checker verifies the required
axiom cases and MP records, with their actual serialization.

The resulting complete local certificate has at most
\begin{equation}\label{eq:local-bound}
 20+(40M+100)Q_b(M)
\end{equation}
characters. The theorem substituting a source-file budget $k$ explicitly
requires $M\le k$. The complete source lexer and its connection to the
agenda's file coding have not been verified.

These local equalities express the definitions' reflexive correctness after
expansion. They do not yet assert the original file's acceptance by the fixed
PA-bin $\operatorname{Bew}$ predicate.

\subsection{A single coded witness for the arithmetic trace}
The next milestone packages all summaries into one natural number, rather
than leaving an external list of semantic values. Use the injective pairing
function
\[
 \pi(a,b)=(a+b)^2+a+1.
\]
For $a+b=z$, its values occupy the interval
$[z^2+1,z^2+z+1]$; these intervals are disjoint, and the value determines $a$
within the interval. A bounded-search decoder and its inverse property are
proved in Lean. Encode a triple by $\pi(L,\pi(S,V))$, and a list by zero for
nil and pairing for cons.

The arithmetic-trace checker validates all instructions against decoded
register values. Lean proves soundness, completeness for the intended trace,
and rejection of missing entries. It may accept irrelevant extra entries;
soundness identifies every register required by the program. A further
compiler builds a compact arithmetic term for the entire encoded trace, and
proves both acceptance and correct lookup of every source summary.

For $n$ source definitions, trace packing adds at most
\begin{equation}\label{eq:trace-bound}
 (8n+20)(112n)
\end{equation}
characters. Since $n\le M$, the total bound is
\begin{equation}\label{eq:bundle-bound}
 20+(40M+100)Q_b(M)+(8M+20)(112M),
\end{equation}
still cubic for fixed $b$. The bounded-search decoder is not claimed to run
in polynomial time; its huge input value may itself be represented compactly.

\subsection{Checked examples and their interpretation}
A doubling grammar with 64 doublings has mass $194$ and expands to a word of
length $2^{64}$. The generated local certificate has 49,287 characters, and
the trace-packing extension has 55,677, for a total of 104,964. These are
generated syntax sizes, not minimum PA proof lengths. The large code is not
evaluated by the demo.

A small boundary-crossing example expands to 13 characters, has mass $14$,
and has a complete generated bundle of 5,443 characters. Independent Python
diagnostics check its printed local file and numeric trace round trip; those
diagnostics are tests, not Lean kernel proofs of the full PA-bin parser.

\status{This milestone verifies the parsed-grammar compiler, written-size
bounds, local restricted proofs, concrete arithmetic encoding, and typed trace
checker. It does not yet give a literal PA formula for the expansion relation,
a polynomial-size PA proof of trace correctness, full Enderton axiom checking,
or agreement with the actual fixed $\operatorname{Bew}$. The remaining internal
PA proof-cost step cannot be replaced by the fact that Lean proves semantic
correctness.}

\section{Proof assembly and the finite-diagonal obstruction}
\label{sec:obstruction}

The next result explains why short ordinary proofs of provability, by
themselves, work against the desired witness. It accounts for the written cost
of references when appending a plain proof suffix to a compressed prefix.

\subsection{A concrete suffix with charged binary references}
Suppose a reflection proof has written length $r$, the target has size $s$,
and an independent \emph{plain} proof of $\Box P$ has length $b_0$. Define
\[
 L(x)=\lfloor\log_2(x+2)\rfloor+1,\qquad
 a=(b_0+s+7)\bigl(2L(b_0+1)+5\bigr).
\]
The suffix construction relocates the plain proof's line references and
appends the final MP step using the reflection conclusion. Plainness avoids
collisions between the suffix's abbreviation definitions and the prefix's
names. With the structured-file validity interface, Lean proves validity and
the serialized bound
\begin{equation}\label{eq:assembly}
 d\le r+aL(r).
\end{equation}
The prefix is represented through its expanded formula history; axioms,
implication, and generalization are parameterized. Thus the record encoding
and reference cost are concrete, but this is not a complete PA-bin parser.

\begin{proposition}[Necessary cost of separation]\label{prop:cost}
Let $q\ge1$ and $t\ge0$. If \eqref{eq:assembly} holds and
$(q+1)r+qt\le qd$, then
\begin{align}
 r+qt&\le qaL(r),\label{eq:cost-necessary}\\
 r&\le8(qa)^2+1,\label{eq:r-obstruction}\\
 d&\le(a+1)\bigl(8(qa)^2+4\bigr).\label{eq:d-obstruction}
\end{align}
\end{proposition}
\begin{proof}
Substitute \eqref{eq:assembly} into the separation inequality and cancel
$qr$ to obtain \eqref{eq:cost-necessary}. The elementary estimate
$L(r)^2\le4(r+2)$, together with $r\le qaL(r)$, gives
$r^2\le4(qa)^2(r+2)$. If $qa=0$, then $r=0$. Otherwise, for $r\ge2$
we have $r+2\le2r$, hence $r\le8(qa)^2$; the small cases are absorbed
by the extra $1$ in \eqref{eq:r-obstruction}. Finally $L(r)\le r+3$ gives
$d\le r+a(r+3)\le(a+1)(r+3)$, proving \eqref{eq:d-obstruction}.
\end{proof}

\begin{corollary}[An eventual exclusion threshold]\label{cor:exclude}
Suppose a family satisfies
\[
 d_n>2^n,\qquad d_n\le r_n+a_nL(r_n),\qquad
 a_n\le A(n+1)^e.
\]
For fixed $q\ge1$, set
\[
 C=(A+1)\bigl(8(qA)^2+4\bigr).
\]
Then factor $1+1/q$ separation, including any nonnegative toll, is impossible
at every index
\begin{equation}\label{eq:threshold}
 n>8(C+3e)^2+1.
\end{equation}
\end{corollary}
\begin{proof}
If separation held, \eqref{eq:d-obstruction} and the bound on $a_n$ would give
$d_n\le C(n+1)^{3e}$. The elementary growth estimate formalized in
\texttt{polynomial\_lt\_exponential} gives
$C(n+1)^{3e}<2^n$ under \eqref{eq:threshold}, a contradiction. The
constants are deliberately coarse; the result excludes every sufficiently
late index, not just some subsequence.
\end{proof}

\subsection{Why the original finite G\"odel family triggers the issue}
For the ordinary proof predicate, let
\[
 B_n=\Box_{2^n}G_n,\qquad \PA\vdash G_n\leftrightarrow\neg B_n.
\]
Inside PA, the two cases already imply $\Box G_n$:
\begin{enumerate}
\item $B_n\to\Box G_n$ by forgetting the length bound.
\item $\neg B_n\to\Box\neg B_n$ by formalized $\Sigma_1$ completeness
      for a $\Sigma_1$ presentation of the decidable bounded-search condition.
\item Necessitating $\neg B_n\to G_n$ and applying distribution gives
      $\Box\neg B_n\to\Box G_n$.
\end{enumerate}
No consistency assumption is needed for this internal argument. Consistency
is used externally to prove the hardness of direct proofs. With uniform
arithmetization, this mechanism supplies a uniform theorem asserting the
provability of the instances. Short specializations of that theorem are the
source of the proposed polynomial plain proofs of $\Box G_n$.

Completing those polynomial proof bounds for the fixed PA-bin representation
would activate Corollary~\ref{cor:exclude}; it would not finish a positive
part~1 witness using this family. The Lean theorem
\texttt{finite\_diagonal\_box} checks the logical implication with its three
object-theory inputs explicit. It does not formalize the PA diagonal lemma
or infer one PA theorem with an object-level universal quantifier from a
metatheoretic family of instance proofs.

\subsection{A stronger-consistency family and the remaining gap}
Consider instead
\[
 T=\PA+\Con(\PA),\qquad Q_n=\Con(T)\restriction n,
\]
using a fixed ordinary checker. If $T$ is consistent, each standard finite
consistency instance $Q_n$ has a PA proof by finite checking. However,
\begin{equation}\label{eq:no-uniform-box}
 \PA\nvdash\forall n\,\Box_{\PA}Q_n.
\end{equation}
Here the box uses the usual represented numeral-instance function.

To prove \eqref{eq:no-uniform-box}, suppose PA proved the displayed universal
statement. Inside $T$, decidable completeness gives
$\neg Q_n\to\Box_{\PA}\neg Q_n$. Together with $\Box_{\PA}Q_n$, this
would imply $\Box_{\PA}\bot$, contrary to $T$'s axiom $\Con(\PA)$.
Thus $T$ would prove $\forall n\,Q_n$, namely $\Con(T)$, contradicting the
second incompleteness theorem for consistent $T$.

This excludes the particular uniform theorem that powered the first family;
it does not exclude all externally obtained short provability proofs. A
suitable short PA proof of $\Box_{\PA}Q_n\to Q_n$ is still missing, as is
the matching lower bound for the exact compressed calculus. A reflection
proof in $T$ about the PA box would mix the two theories and would not answer
the question. The slow-consistency variant does not automatically repair the
route: Freund and Pakhomov obtain uniform PA-provability of the corresponding
finite instances in their setting \cite[Corollary 2.4; Proposition 3.7]{fp}.

\subsection{Why no universal cheap-box shortcut is available}
Under $\Sigma_1$ soundness of ordinary PA, there is no total computable
function $g$ bounding $\ell_{\PA}(\Box P)$ by $g(|P|)$ for every PA theorem
$P$. Otherwise, on input $P$, compute the bound and inspect all PA files up
to that length for a proof of $\Box P$. If $P$ is a theorem, the bound makes
the search succeed; if the search succeeds, $\Sigma_1$ soundness of the proof
predicate makes $P$ provable. This would decide PA theoremhood. This argument
rules out a universal computable bound, not family-specific estimates.

\status{The suffix construction, inequalities, threshold, and parameterized
finite-diagonal logic are Lean-verified. The arithmetical arguments about
$T$, the uniform boxed families, and theoremhood undecidability are
mathematical arguments in this report; they are not complete PA
formalizations in the supplied projects.}

\section{What remains for ordinary PA-bin}
\label{sec:remaining}

The work has identified both a structural target and a concrete implementation
target. For part~1, a positive answer requires non-sublinearity of $H$ for the
actual box, while a negative answer requires $H=o(k)$, equivalently the
near-unit sublinear additive bound of Theorem~\ref{thm:envelope}. Merely
showing some polynomial conversion cost chooses neither alternative.

For the polynomial route in part~2, the remaining internalization chain is:
\begin{enumerate}
\item Connect literal input bytes, identifier decoding, and the source
      G\"odel code to the well-scoped raw-string grammar, with a charged bound
      relating grammar mass to file length.
\item Write the actual PA expansion/trace formula. Generate short PA proofs
      of the arithmetic pairing, sequence-access, and instruction-correctness
      facts needed by the compact witness. Bound the \emph{written object
      proofs}, not only their Lean semantics.
\item Cover the full Enderton checker, including propositional-tautology
      axioms, capture-avoiding substitution, generalization, MP, and the
      abbreviation scope rules.
\item Prove agreement with the agenda's fixed $\operatorname{Bew}$ and
      produce the notice/internal L\"ob operations required by
      \texttt{QuantitativeLobTools}, with explicit character bounds.
\end{enumerate}
These are unresolved obligations, not established lemmas to be imported on
trust. The raw-string quotation work handles a significant representation
obstacle, while the suffix obstruction explains why completing short-box
proofs for the first hard family would eliminate that candidate rather than
create the required separation.

\section{Certificate ledger and local reproduction}
\label{sec:ledger}

\subsection{What the cumulative archive contains}
All six projects are pinned to \textbf{Lean 4.19.0}. They use the libraries
bundled with Lean and need no Mathlib checkout. The archive has this structure:
\begin{lstlisting}
MAIS-O11-certificates/
  verify.sh
  README.md
  SHA256SUMS
  verified-build.txt
  projects/
    MAISO11FirstPass/
    MAISO11AuditPass/
    MAISO11TaggedPass/
    MAISO11QuotationPass/
    MAISO11ObstructionPass/
    MAISO11EnvelopePass/
  research-notes/
\end{lstlisting}
The research notes preserve the detailed milestone accounts and the original
slow-verifier manuscript. Read their claims with the current scope distinctions
in this report. The projects retain their own original verification transcripts;
\texttt{verified-build.txt} records the cumulative runner's new build.

{\small
\begin{longtable}{@{}>{\raggedright\arraybackslash}p{0.25\textwidth}>{\raggedright\arraybackslash}p{0.69\textwidth}@{}}
\toprule
Project & Certified content and principal declarations \\
\midrule
\endfirsthead
\toprule Project & Certified content and principal declarations \\
\midrule
\endhead
\bottomrule\endfoot
\texttt{FirstPass} & Conditional finite diagonal contradiction, abstract
L\"ob rule, lower-bound transfer, shortest proofs, and divergence.
\path{MAISO11.SeparationHypotheses.conditional_main}. \\[5pt]
\texttt{AuditPass} & Internal L\"ob logic and polynomial conversion from
explicit quantitative operations.
\path{MAISO11.internal_lob_axiom},
\path{MAISO11.QuantitativeLobTools.uniform_budget},
\path{MAISO11.QuantitativeLobTools.linear_premise_bound},
\path{MAISO11.first_pass_and_polynomial_tools_incompatible}. \\[5pt]
\texttt{TaggedPass} & Concrete tagged syntax and guard preservation, expanded
trace acceptance under stated inputs, and numerical gap implications.
\path{MAISO11.Tagged.shape_subst},
\path{MAISO11.Tagged.guarded_instantiation_axiom},
\path{MAISO11.Tagged.no_polynomial_bridges}. \\[5pt]
\texttt{QuotationPass} & Raw-string compiler, local certificates, arithmetic
trace checker, compact witness, and serialized cost. In namespace
\path{MAISO11.Quotation}:
\path{compile_correct}, \path{compile_certificate_characters},
\path{compiled_localCertificate_checked}, \path{compiled_trace_sound},
\path{packedWitness_checked}, \path{compiled_bundle_characters}. \\[5pt]
\texttt{ObstructionPass} & Valid serialized suffix and explicit cost;
separation obstruction and eventual threshold; abstract finite-diagonal
logic. In namespace \path{MAISO11.Obstruction}:
\path{Suffix.appendBoxSuffix_normalized},
\path{separation_reflection_bound}, \path{no_separation_after},
\path{finite_diagonal_box}. \\[5pt]
\texttt{EnvelopePass} & Exact negative condition, affine bounds, and arbitrary
factor consequence. In namespace \path{MAISO11.Envelope.Lengths}:
\path{noSpeedup_iff_sublinear_envelope},
\path{noSpeedup_iff_sublinear_converter},
\path{noSpeedup_near_unit_affine},
\path{noSpeedup_overhead_affine},
\path{not_linear_overhead_arbitrary_factor}.
Finite files: \path{MAISO11.Envelope.FiniteFiles.toLengths}. \\
\end{longtable}
}

\subsection{Run the certificates on a Mac}
Download \texttt{MAIS-O11-certificates.zip} to Downloads. The following uses
a fresh temporary folder, so it also works when repeating the verification:
\begin{lstlisting}[language=bash]
mais_check_dir="$(mktemp -d "${TMPDIR:-/tmp}/mais-o11.XXXXXX")"
unzip -q "$HOME/Downloads/MAIS-O11-certificates.zip" -d "$mais_check_dir"
cd "$mais_check_dir/MAIS-O11-certificates"
bash verify.sh
\end{lstlisting}
The runner locates \texttt{lake}, including the usual
\texttt{\$HOME/.elan/bin/lake} fallback, and rebuilds every project. It then
runs each \texttt{Audit.lean} with warnings treated as errors. The toolchain
file selects Lean~4.19.0; Elan may download that release if it is not already
installed. The runner prints a final success message only if all six projects
pass. It rebuilds the proof terms; it does not just print a saved transcript.

To recheck only the latest result from the archive root:
\begin{lstlisting}[language=bash]
cd projects/MAISO11EnvelopePass
lake build
lake env lean -DwarningAsError=true Audit.lean
\end{lstlisting}
The quotation demo is optional and is not needed for the theorem checks:
\begin{lstlisting}[language=bash]
cd projects/MAISO11QuotationPass
lake exe quotationdemo
\end{lstlisting}
Run that last block from the archive root, rather than from inside the envelope
project. The independent Python diagnostics described in the quotation
project's README are separate tests of printed examples.

\subsection{How to interpret the audit}
The audited theorem dependencies are contained in Lean's standard logical
axioms \texttt{propext}, \texttt{Quot.sound}, and
\texttt{Classical.choice}; some theorems need fewer or none. The supplied
Lean source has no admitted proofs, custom axiom declarations,
\texttt{native\_decide} proofs, or unsafe declarations.

An axiom report is \emph{not} a list of a theorem's hypotheses. A conditional
theorem may have no nonstandard axioms while still taking the desired
arithmetization or proof-cost estimate as a parameter. Read its printed type
and the interface definitions. Moreover, lengths in these theorems refer to
the specified object certificates; they do not count Lean source characters,
tactic steps, or compiled \texttt{.olean} bytes.

\subsection{Compile this report}
The TeX source is self-contained, with its bibliography included. A standard
TeX Live or MacTeX installation can produce the matching PDF using:
\begin{lstlisting}[language=bash]
pdflatex -halt-on-error MAIS-O11-progress-report.tex
pdflatex -halt-on-error MAIS-O11-progress-report.tex
\end{lstlisting}
Two passes resolve the table of contents and cross-references. The PDF and
TeX contain the same consolidated account; the companion archive supplies
the executable certificates and detailed source notes.

\clearpage
\phantomsection
\addcontentsline{toc}{section}{References}
\begin{thebibliography}{9}
\bibitem{mais}
MAIS project. \emph{Does L\"ob's rule shorten proofs?} (MAIS-O11),
\emph{MAIS-A1}, Sections 2--3 and Question 3.6, and MAIS-O10.
Uploaded source documents used in this development; copies are included in
the research-notes directory.
\url{https://github.com/lionellevine/MAIS}.

\bibitem{enderton}
Herbert B. Enderton. \emph{A Mathematical Introduction to Logic},
second edition. Academic Press, 2001. Section 2.4.

\bibitem{critch}
Andrew Critch. A parametric, resource-bounded generalization of L\"ob's
theorem, and a robust cooperation criterion for open-source game theory.
\emph{Journal of Symbolic Logic} 84(4), 1368--1381, 2019.
\url{https://arxiv.org/abs/1602.04184}.

\bibitem{pudlak}
Pavel Pudl\'ak. The lengths of proofs. In Samuel R. Buss, editor,
\emph{Handbook of Proof Theory}, pp. 547--637. Elsevier, 1998.
\url{https://users.math.cas.cz/~pudlak/length.pdf}.

\bibitem{finite}
Pavel Pudl\'ak. Incompleteness in the finite domain.
\emph{Bulletin of Symbolic Logic} 23(4), 405--441, 2017.
\url{https://arxiv.org/abs/1601.01487}.

\bibitem{fp}
Anton Freund and Fedor Pakhomov. Short proofs for slow consistency.
\emph{Notre Dame Journal of Formal Logic} 61(1), 31--49, 2020.
\url{https://arxiv.org/abs/1712.03251}.
\end{thebibliography}
\end{document}
